\documentclass[10pt,letterpaper]{article}
\usepackage[T1]{fontenc}
\usepackage[utf8]{inputenc}
\usepackage{lmodern}
\usepackage[letterpaper,margin=0.84in,headheight=15pt]{geometry}
\usepackage[table]{xcolor}
\definecolor{DeepNavy}{HTML}{173450}
\definecolor{Teal}{HTML}{087F8C}
\definecolor{Slate}{HTML}{536578}
\definecolor{SoftBlue}{HTML}{F1F6FA}
\definecolor{SoftGray}{HTML}{F6F8FA}
\definecolor{Crimson}{HTML}{B34843}
\usepackage{microtype}
\setlength{\emergencystretch}{1.5em}
\usepackage{booktabs,longtable,tabularx,array}
\usepackage{amssymb}
\usepackage{ragged2e}
\usepackage{enumitem}
\usepackage{fancyhdr,lastpage}
\usepackage{titlesec}
\usepackage{listings}
\usepackage[most]{tcolorbox}
\usepackage{needspace}
\usepackage{hyperref}
\hypersetup{colorlinks=true,linkcolor=DeepNavy,urlcolor=Teal,citecolor=DeepNavy,pdfauthor={Software Engineering Technical Documentation}}
\setlength{\parindent}{0pt}
\setlength{\parskip}{5pt plus 1pt}
\setlist[itemize]{itemsep=2pt,topsep=3pt,leftmargin=1.55em}
\setlist[enumerate]{itemsep=2pt,topsep=3pt,leftmargin=1.7em}
\titleformat{\section}{\Large\bfseries\color{DeepNavy}}{\thesection}{0.75em}{}[\vspace{-3pt}\color{Teal}\titlerule]
\titleformat{\subsection}{\large\bfseries\color{DeepNavy}}{\thesubsection}{0.65em}{}
\titleformat{\subsubsection}{\normalsize\bfseries\color{Teal}}{\thesubsubsection}{0.55em}{}
\pagestyle{fancy}
\fancyhf{}
\fancyhead[L]{\sffamily\footnotesize MISRA C:2023 | Engineering Documentation}
\fancyhead[R]{\sffamily\footnotesize\color{Slate} Independent companion}
\fancyfoot[L]{\sffamily\scriptsize Original explanatory material; consult licensed MISRA standard}
\fancyfoot[R]{\sffamily\scriptsize\thepage\ / \pageref{LastPage}}
\renewcommand{\headrulewidth}{0.3pt}
\newcommand{\ID}[1]{\textsf{\bfseries #1}}
\newcommand{\SourceNote}[1]{\par\smallskip{\footnotesize\color{Slate}\textit{Trace: #1.}}\par}
\newtcolorbox{keybox}[1][]{enhanced,breakable,colback=SoftBlue,colframe=Teal,boxrule=0.5pt,arc=1mm,title={Key point},fonttitle=\bfseries,#1}
\newtcolorbox{riskbox}[1][]{enhanced,breakable,colback=SoftGray,colframe=Crimson,boxrule=0.5pt,arc=1mm,title={Verification boundary},fonttitle=\bfseries,#1}
\lstdefinestyle{engineerc}{language=C,basicstyle=\ttfamily\footnotesize,keywordstyle=\color{DeepNavy}\bfseries,commentstyle=\color{Slate}\itshape,stringstyle=\color{Teal},numbers=left,numberstyle=\tiny\color{Slate},stepnumber=1,numbersep=8pt,breaklines=true,breakatwhitespace=false,frame=single,framesep=5pt,rulecolor=\color{Slate!25},backgroundcolor=\color{SoftGray},tabsize=2,showstringspaces=false,columns=fullflexible,keepspaces=true}
\lstdefinestyle{plaincode}{basicstyle=\ttfamily\footnotesize,numbers=none,breaklines=true,frame=single,rulecolor=\color{Slate!25},backgroundcolor=\color{SoftGray},columns=fullflexible,keepspaces=true}
\newcommand{\TrainingTag}[1]{\textsf{\footnotesize\color{Teal}\bfseries #1}}
\newcommand{\DocNotes}{\begin{keybox}This publication is an independently written educational and operational companion. It is \emph{not} the licensed MISRA guideline text, a MISRA-endorsed tool, or a certification of project compliance. Identifiers and published classifications are used for traceability; check exact requirement wording, amplification, applicability, and exceptions in the licensed standard.\end{keybox}}
\AtBeginDocument{\color{black}}

\title{\vspace{-2.3em}\color{DeepNavy}\Huge\bfseries Compliance Operating Procedure\\[0.45em]\large\normalfont\color{Teal} Lifecycle controls, GEP, GRP, GCS and release evidence}
\author{Independent Software Engineering Reference Series}
\date{Edition 1.0 \quad | \quad October 2026}
\begin{document}
\maketitle
\vspace{-0.4em}
\begin{center}\color{Slate}\rule{0.83\linewidth}{0.4pt}\end{center}
\DocNotes
\tableofcontents
\newpage

\section{Compliance means a project claim}
A claim needs explicit scope: repository revision, deliverable and binary, target platform, C edition, vendor components, generated code and checked build configurations. A company-level claim or a claim based solely on green static analysis cannot replace the project-specific compliance process defined in MISRA Compliance:2020.
\section{Roles and accountability}
\begin{tabularx}{\linewidth}{@{}p{.28\linewidth}X@{}}\toprule
Role & Accountability\\\midrule
Project owner / acquirer & Accepts applicable standards, scope, re-categorization and residual risk.\\
Engineering lead & Maintains source inventory, compiler configurations, traceability and code review.\\
Analyzer owner & Maintains tools, guideline enforcement mapping, validation tests and report reproducibility.\\
Assurance / safety / AppSec reviewer & Challenges coverage, deviations and risk claims independently.\\
Approver & Accepts specifically allowed deviations and signs final compliance evidence.\\\bottomrule\end{tabularx}
\section{Lifecycle stage gates}
\subsection{Gate 0: initiate}
Select source revision and language edition; declare scope, vendor interfaces, risk classification, functional requirements, tooling and escalation owner. Create a \textbf{Guideline Enforcement Plan (GEP)} identifying checking methods for all applicable entries.
\subsection{Gate 1: design and build}
Determine implementation-defined behaviors (Dir 1.1), code-to-requirement traceability (Dir 3.1), and runtime checks (Dir 4.1). Build with no compiler errors. Ensure test cases exist for negative and boundary conditions.
\subsection{Gate 2: verification}
Run the documented GEP methods with frozen configurations. Record diagnostics and verify coverage of multi-TU and undecidable items. Maintain a \textbf{Guideline Re-categorization Plan (GRP)} if project-specific category strengthening is agreed.
\subsection{Gate 3: disposition}
Repair violations or route through the formal deviation process only where permitted. Track unresolved diagnostic uncertainty, tool gaps and supplier evidence as issues, not as silent exclusions.
\subsection{Gate 4: release approval}
Prepare a \textbf{Guideline Compliance Summary (GCS)} with one status per relevant guideline, backed by documented evidence; include deviation records and any permits, signoffs and delivery boundaries.
\section{Status accounting}
For a project-compliance summary, MISRA Compliance:2020 uses \emph{Compliant}, \emph{Deviations}, \emph{Violations}, and \emph{Disapplied}. Their permitted combinations depend on category and any agreed re-categorization. This companion does not redefine those status rules; use the compliance publication to validate every final GCS row.
\begin{riskbox}Marking a guideline ``Disapplied'' is \emph{not} equivalent to ``not applicable''. The former concerns intentionally not enforcing a guideline under a permitted compliance arrangement; the latter depends on the rule's edition applicability or project-specific exclusion and must be justified.\end{riskbox}
\section{Supplier and generated code treatment}
Inventory standard library, target vendor drivers, operating-system components, generated source and binary-only modules. Attach supplier statements where available, verify boundary headers and assess integration risk independently. For automatically generated code inspect the publication's Appendix E; categories may differ from hand-authored code and injected user code cannot automatically claim generator treatment.
\section{Evidence structure}
\begin{lstlisting}[style=plaincode]
compliance/
  scope-manifest.yml
  guideline-enforcement-plan.csv
  guideline-recategorization-plan.csv
  guideline-compliance-summary.csv
  diagnostics/   # immutable analysis snapshots
  deviations/    # approved risk decisions
  tools/         # versions, digests, tests, configs
  requirements/  # code-to-requirement links
  suppliers/     # claims and assessment of gaps
  release/       # final audit and sign-offs
\end{lstlisting}
\section{Audit sampling protocol}
Sample at least one guideline from each topic, all project-specific mandatory items, all approved deviations, system-level analyses, resource lifetime checks and any supplier-provided modules. For each selected item reproduce the checker invocation and trace the outcome to the final GCS. Define the actual sampling plan according to project integrity requirements.
\section{Traceability chain}
\begin{center}\fbox{Requirement $\rightarrow$ Design $\rightarrow$ Source $\rightarrow$ Check/Test $\rightarrow$ Finding $\rightarrow$ Fix/Deviation $\rightarrow$ GCS}\end{center}
\SourceNote{MISRA C:2023 Sections 1.4, 6--8 and Appendices A, B, E, I; MISRA Compliance:2020}

\end{document}
